\section{本章小结}

Shell 的核心数据结构：

\begin{enumerate}
  \item \textbf{\texttt{cmd\_t}}：name + help + fn 函数指针，每个命令一个独立编译单元
  \item \textbf{\texttt{g\_cmds[]}}：静态指针数组，编译期确定，遍历分发
  \item \textbf{行缓冲}：栈上 128 字节，tokenize 原地修改（零拷贝）
  \item \textbf{Console 层}：统一键盘 + 串口输入，屏蔽输入源差异
\end{enumerate}

Shell 是内核开发阶段最重要的调试工具——后续每实现一个子系统
都会立即添加对应的调试命令。下一章进入内存管理：PMM（物理内存分配器）。

